\documentclass[11pt]{article}

\usepackage[a4paper,margin=2.4cm]{geometry}
\usepackage{amsmath,amssymb}
\usepackage{booktabs}
\usepackage{siunitx}
\usepackage{xcolor}
\usepackage[hidelinks]{hyperref}
\usepackage{enumitem}

\newcommand{\bplus}{\beta^{+}}
\newcommand{\Cl}[1]{{}^{#1}\mathrm{C}}
\newcommand{\Ox}[1]{{}^{#1}\mathrm{O}}
\newcommand{\Nt}[1]{{}^{#1}\mathrm{N}}

\title{A simulation chain to compare PET detectors for\\
proton-therapy range verification}
\author{Plan / working document}
\date{\today}

\begin{document}
\maketitle

\begin{abstract}
We describe a simulation machinery whose goal is \emph{not} to reproduce the
clinical study of Parodi, Bortfeld \& Haberer (\textit{Int.\ J.\ Radiat.\ Oncol.\
Biol.\ Phys.}\ \textbf{71}, 945, 2008), but to use its physics inputs as a
realistic context in which to \textbf{compare candidate PET detectors} for
verification of the proton range in the \textbf{in-room} modality. The central
design principle is to \textbf{decouple production from detection}: a single,
detector-independent event file is generated once with Geant4 and then passed
through the response of each candidate PET. Differences in the reconstructed
images are therefore attributable to the detectors alone, which makes the
comparison robust against the (imperfect) nuclear cross sections of the transport
code. The figure of merit is the precision with which the \emph{distal fall-off}
of the activity profile -- and hence the proton range -- is recovered at the
photon-starved counting statistics typical of this measurement.
\end{abstract}

\section{Overview of the chain}
\label{sec:overview}

The pipeline is two Geant4 stages joined by a file, followed by an independent
reconstruction stage. Stage~A (proton transport, the expensive part) runs
\textbf{once} and emits a detector-independent source file. A short analytic step
then sets how many decays are measured. Stage~B (the detector) runs \textbf{once
per candidate PET} and emits a \textbf{coincidence list}. Reconstruction is a
separate downstream stage that consumes only that list, so it can be swapped or
deferred without touching the simulation.

\begin{center}
\fbox{\parbox{0.95\textwidth}{\ttfamily\small
\textbf{[A] GEANT4 --- transport} (gun + geometry + physics) \hfill \textit{runs ONCE}\\
\hspace*{2em}proton beam $\to$ transport in PMMA $\to$ $\bplus$ emitter produced (residual nucleus)\\
\hspace*{4em}$\to$ $\bplus$ decay (natural lifetime) $\to$ annihilation\\
\hspace*{2em}$\Rightarrow$ \textbf{PROD/ANH file}: per event, per isotope $j$ --- PROD $(x,y,z)$, ANH $(x',y',z')$\\
\\
\textbf{[B$_0$] TIME-DECAY BOOKKEEPING} (analytic, the A$\to$B handoff)\\
\hspace*{2em}production $P_j$ $\to$ $N_j$ decays in the in-room window\\
\hspace*{2em}$\to$ draw $N_j$ ANH events per isotope $\Rightarrow$ \textbf{sampled annihilation events}\\
\\
\textbf{[B] GEANT4 --- detector} ($\sigma_E,\sigma_t$, DOI, attenuation) \hfill \textit{runs PER detector}\\
\hspace*{2em}annihilation events $\to$ detector response $\to$ coincidence sorting\\
\hspace*{2em}$\Rightarrow$ \textbf{COINCIDENCE list} (two DOI 3-D hits, energies, times)\\
\\
\hspace*{2em}-------- independent boundary --------\\
\\
\textbf{[C] RECONSTRUCTION} (separate stage, toolkit deferred)\\
\hspace*{2em}coincidence list $\to$ image $\to$ $\sigma(\text{range})$ vs.\ counts $\to$ compare
}}
\end{center}

\noindent
The key consequences:
\begin{itemize}[nosep]
\item \textbf{Stage~A runs once.} The expensive hadronic proton transport is not
repeated when the detector changes.
\item \textbf{Every candidate PET sees the identical source.} The comparison of
Sec.~\ref{sec:compare} is then free of per-run statistical drift between detectors.
\item \textbf{The $N_j$ sampling lives at the A$\to$B handoff.} A single monolithic
pass would emit one annihilation per \emph{produced} isotope
(production-proportional), not the windowed \emph{measured} number; sampling $N_j$
per isotope there is the only place the $\Ox{15}/\Cl{11}$ mix is set correctly.
\item \textbf{Reconstruction is decoupled.} The simulation's deliverable is the
coincidence list; any reconstructor consumes it.
\end{itemize}
A high-resolution but low-sensitivity PET can lose to a coarser, high-sensitivity
one at the low counts characteristic of in-room acquisition ($\sim\!10^{5}$--$10^{6}$
coincidences for a \SI{1}{Gy} field), and the machinery is built to expose exactly
that trade-off.

\section{Geant4 stage A: transport --- beam, phantom, physics, and the event file}
\label{sec:g4}

This stage runs \textbf{once}, ahead of all detector studies. It requires only
Geant4 ingredients -- a particle gun, a physics list and
a geometry -- and produces a clean, detector-independent file that is the input
to everything downstream. The primary launched by the gun is the \textbf{proton};
the positron emitters are \emph{products} of its nuclear interactions in the
phantom.

\subsection{Beam (the gun)}
\label{sec:beam}
We adopt mainstream scanned-proton parameters. The clean baseline is a single
mono-energetic pencil beam, which produces a one-dimensional depth--activity
profile that is the ideal range-verification testbed; a spread-out configuration
is a later extension.

\begin{table}[h]
\centering
\caption{Baseline beam. Ranges are approximate continuous-slowing-down values in
water. The macroscopic time structure is \emph{not} used by the gun; it enters
only the time-decay bookkeeping of Sec.~\ref{sec:timedecay} (Table~\ref{tab:time}).}
\label{tab:beam}
\begin{tabular}{ll}
\toprule
Species          & proton (${}^{1}$H) \\
Energy (baseline)& \SI{100}{MeV} ($R\approx\SI{7.7}{cm}$); cross-check \SI{150}{MeV} ($R\approx\SI{15.8}{cm}$) \\
Lateral profile  & Gaussian pencil, $\sigma\approx\SI{3}{}$--$\SI{5}{mm}$ at entrance \\
Direction        & along the phantom (cylinder) axis \\
Fluence          & normalised to \SI{1}{Gy} at the target, \emph{or} fixed $N_p$ with dose reported \\
\bottomrule
\end{tabular}
\end{table}

\subsection{Phantom (the geometry)}
\label{sec:phantom}
A right circular cylinder of homogeneous, constant-density material, axis along
the beam. A homogeneous phantom is deliberate: it provides \textbf{ground truth}
and isolates the detector from tissue-heterogeneity and washout effects.

\paragraph{Material: PMMA as baseline.}
We choose \textbf{PMMA} ($\mathrm{C_5H_8O_2}$, $\rho=\SI{1.18}{g/cm^3}$) because
(i) it contains both carbon and oxygen, producing the clinically relevant
$\Cl{11}+\Ox{15}$ mixture -- two isotopes with very different half-lives
\emph{and} positron ranges, which is exactly what stresses a detector; (ii) it is
the de-facto standard solid phantom for proton-PET benchmarking; (iii) it is
reproducible. \textbf{Water} is kept as an optional cross-check: with no carbon it
yields an almost single-isotope ($\Ox{15}$) source, a clean limiting case for the
longest positron range.

\paragraph{Geometry (suggested).}
Diameter \SI{20}{cm}, length \SI{20}{cm} (enough to stop the baseline beams and to
act as a realistic attenuator of the \SI{511}{keV} photons). The cylinder must be
present as a passive attenuator in \emph{every} detector geometry of
Sec.~\ref{sec:detector}, so that attenuation and scatter are treated consistently.

\subsection{Physics list}
\label{sec:physics}
A hadronic list suited to proton-induced isotope production, e.g.\
\texttt{QGSP\_BIC\_HP} (or \texttt{QGSP\_BERT\_HP}), with electromagnetic physics
and radioactive decay enabled (the latter at natural lifetime, see
Sec.~\ref{sec:production}).
\emph{Caveat:} Geant4's isotope-production cross sections are less tuned than
data-driven approaches (FLUKA, or analytic models built on EXFOR). This biases the
\emph{absolute} yields and somewhat the production-profile \emph{shape} -- but,
per Sec.~\ref{sec:compare}, it is common-mode across detectors and cancels in the
ranking. The file can later be regenerated from a FLUKA or measured map with no
change to the downstream blocks.

\subsection{Production of the \texorpdfstring{$\bplus$}{beta+} emitters and the event file}
\label{sec:production}
As each proton slows down in the PMMA it loses energy electromagnetically and,
with a probability set by the inelastic nuclear cross sections, undergoes a
nuclear reaction with a $\Cl{12}$ or $\Ox{16}$ nucleus of the medium. Such a
reaction leaves a \textbf{residual nucleus}; in a fraction of cases that residual
is one of the $\bplus$ emitters of Table~\ref{tab:iso} (e.g.\
$\Cl{12}(p,pn)\Cl{11}$, $\Ox{16}(p,pn)\Ox{15}$), created essentially at rest at
the reaction point.

We let the emitter decay through Geant4's \textbf{standard radioactive-decay
process at its natural lifetime} -- no prompt-decay override, which would be an
unnecessary intervention. Geant4 emits the positron with the correct $\bplus$
spectrum of the species; the neutrino escapes and the (stable) daughter recoils
only $\sim\!\si{\micro m}$. The positron is transported through its
(sub-millimetre) range and annihilates. Geant4 decays \emph{every} produced
emitter exactly once, so the file is \textbf{production-proportional} (one
annihilation per emitter, 1:1); the measured number $N_j$ and the acquisition
timing are set only later, by the bookkeeping of Sec.~\ref{sec:timedecay}. For
each such emitter we write \textbf{two records}:
\begin{center}
\begin{tabular}{llll}
\texttt{mode=PROD} & \texttt{isotope=}$j$ & $x,y,z$    & $t=0$  \\
\texttt{mode=ANH}  & \texttt{isotope=}$j$ & $x',y',z'$ & $t'$   \\
\end{tabular}
\end{center}
The PROD point is the production/decay point (the recoil moved $\sim\!\si{\micro m}$);
the ANH point is where the two \SI{511}{keV} photons are born. Three properties of
this file matter:
\begin{itemize}[nosep]
\item \textbf{PROD is the truth} activity map $P_j(\mathbf{r})$, the reference
against which a reconstruction's range is judged.
\item \textbf{ANH is the detector source}, fed to Sec.~\ref{sec:detector}.
\item \textbf{The displacement} $|\text{ANH}-\text{PROD}|$ \textbf{is the positron
range}, per event -- the resolution floor comes out of the same physics, with no
separate kernel needed.
\end{itemize}

\paragraph{On the time stamps.}
With natural decay, $t'$ is the physically sampled decay time (seconds to
minutes) plus the positron's $\sim\!\si{ps}$ flight time. It carries no
acquisition-time information for this study and is \textbf{not stored} -- the
event file records only the two space points. This is by design: \textbf{Geant4
owns space (the two points); the time-decay bookkeeping of
Sec.~\ref{sec:timedecay} owns number and acquisition timing.} Because an isotope
does not move between production and decay, \emph{where} it annihilates is
independent of \emph{when} it decays, so the spatial file is complete on its own
and the decay time can be dropped.

\paragraph{Backgrounds discarded here (small, noted).}
$\Cl{10}$ and $\Ox{14}$ emit a prompt de-excitation $\gamma$ in coincidence with
the $\bplus$; killing all but the positron removes it. It is a minor, genuine
background, and these isotopes are largely gone at the in-room delay, so it is
acceptable for a first cut. Prompt neutrons from the proton shower can cause
secondary activation; killing them for speed loses a small contribution to
$P_j(\mathbf{r})$.

\paragraph{Equivalent alternative.}
One may instead write only the PROD record (no decay at all) and add the
positron-range blur downstream by convolving $P_j(\mathbf{r})$ with a precomputed
isotope-specific kernel. This is equivalent but strictly less informative -- the
two-point file already contains the per-event range -- so we adopt the two-point
scheme as baseline.

\subsection{Normalisation to dose}
\label{sec:norm}
Score the deposited dose in the phantom in the same run. Geant4 then yields both
the production per primary, $P_j/N_p$ (PROD records / protons simulated), and the
dose per primary, $D/N_p$. Scaling to a clinical \SI{1}{Gy} field fixes the
physical proton number $N_p(\SI{1}{Gy}) = \SI{1}{Gy}\,/(D/N_p)$ and hence the
integral production $P_j = (P_j/N_p)\,N_p(\SI{1}{Gy})$ that the bookkeeping of
Sec.~\ref{sec:timedecay} consumes. This is the bridge between the (arbitrary)
number of simulated primaries and the absolute count budget.

\section{Time-decay bookkeeping}
\label{sec:timedecay}

This block is analytic. It takes the integral production $P_j$ from
Sec.~\ref{sec:norm} and returns $N_j$, the number of decays of species $j$ that
fall inside the chosen acquisition window for a given beam time structure. Each
species has decay constant $\lambda_j=\ln 2/T_{1/2,j}$ and is assumed to decay
$100\%$ by $\bplus$ emission. The model is that of Parodi et al.\ (their
Eqs.~1--7).

\begin{table}[h]
\centering
\caption{Positron emitters considered, with half-lives and the dominant
production channel in the phantom media. The $\bplus$ endpoint energy governs the
positron range. Values are nominal; adopt tabulated nuclear data in the
implementation.}
\label{tab:iso}
\begin{tabular}{lccl}
\toprule
Isotope & $T_{1/2}$ & $\bplus$ endpoint (MeV) & Dominant channel \\
\midrule
$\Ox{15}$ & \SI{122}{s}   & 1.74 & $\Ox{16}(p,pn)\Ox{15}$ \\
$\Cl{11}$ & \SI{1223}{s}  & 0.96 & $\Cl{12}(p,pn)\Cl{11}$ \\
$\Nt{13}$ & \SI{598}{s}   & 1.19 & on O / N \\
$\Cl{10}$ & \SI{19.3}{s}  & 1.91 & $\Cl{12}(p,pn n)$; \textit{prompt $\gamma$ \SI{718}{keV}} \\
$\Ox{14}$ & \SI{70.6}{s}  & 1.81 & on O; \textit{prompt $\gamma$} \\
\bottomrule
\end{tabular}
\end{table}

\paragraph{Production rate.}
For a delivery of $S$ spills of duration $t_s$, with the constant-activation
approximation, $K_j = P_j/(S\,t_s)$.

\paragraph{Activity within and between spills (synchrotron).}
With extraction period $t_{\mathrm{extr}}=t_s+t_p$ (spill $+$ pause),
\begin{align}
A_j(t) &= K_j\left(1-e^{-\lambda_j t}\right), & 0\le t\le t_s, \\
A_j(t) &= K_j\left(1-e^{-\lambda_j t_s}\right)e^{-\lambda_j (t-t_s)}, & t_s < t < t_{\mathrm{extr}},
\end{align}
and the number of decays in an interval $\Delta t$ starting at $t$ is
$N_j(t,\Delta t)=\int_t^{t+\Delta t} A_j(t')\,dt'$.

\paragraph{In-beam decays (retained for completeness).}
\begin{align}
N_{j,\mathrm{spill}} &= S\,\Delta_{\mathrm{spill}}
  + \frac{K_j}{\lambda_j}\left(1-e^{-\lambda_j t_s}\right)^2 e^{-\lambda_j t_p}
    \sum_{i=1}^{S-1}(S-i)\,e^{-(i-1)\lambda_j t_{\mathrm{extr}}},
    \label{eq:spill}\\[2pt]
\Delta_{\mathrm{spill}} &= K_j t_s - \frac{K_j}{\lambda_j}\left(1-e^{-\lambda_j t_s}\right),\\[4pt]
N_{j,\mathrm{pause}} &= \frac{K_j}{\lambda_j}\left(1-e^{-\lambda_j t_s}\right)\left(1-e^{-\lambda_j t_p}\right)
    \sum_{i=1}^{S-1}(S-i)\,e^{-(i-1)\lambda_j t_{\mathrm{extr}}},
    \label{eq:pause}\\[4pt]
N_{j,\mathrm{end}} &= \frac{K_j}{\lambda_j}\left(1-e^{-\lambda_j t_s}\right)\left(1-e^{-\lambda_j t_{\mathrm{end}}}\right)
    \sum_{i=0}^{S-1} e^{-i\lambda_j t_{\mathrm{extr}}}. \label{eq:end}
\end{align}

\paragraph{The measured (in-room) window.}
A window $t_{\mathrm{meas}}$ starting a delay $t_{\mathrm{del}}$ after the end of
irradiation collects
\begin{equation}
N_{j,\mathrm{meas}} = \frac{K_j}{\lambda_j}\left(1-e^{-\lambda_j t_s}\right) e^{-\lambda_j t_{\mathrm{del}}}\left(1-e^{-\lambda_j t_{\mathrm{meas}}}\right)
    \sum_{i=0}^{S-1} e^{-i\lambda_j t_{\mathrm{extr}}}. \label{eq:off}
\end{equation}

\paragraph{Cyclotron (continuous) limit.}
A continuous delivery is the single-spill case: $S=1$ and
$t_{\mathrm{extr}}=t_s=t_{\mathrm{irr}}$; the recursive sums collapse
($\sum_{i=1}^{0}\!\equiv 0$, $\sum_{i=0}^{0}\!\equiv 1$).

\paragraph{In-room scope.}
This study targets the \textbf{in-room} modality: a stand-alone full-ring scanner
in the treatment room imaging the patient on the couch a short delay after
irradiation, without repositioning. The defining parameter is a \emph{short}
delay; the minimum is $t_{\mathrm{del}}\approx\SI{120}{s}$ (the MGH in-room mean
was $\approx\SI{2.5}{min}$), and larger delays are part of the study via the
dependence on $t_{\mathrm{del}}$. The count budget is Eq.~\eqref{eq:off} at this
short delay. The physical consequence that drives the detector study: at such a
short delay the short-lived $\Ox{15}$ ($T_{1/2}=\SI{122}{s}$) is still abundant
and \emph{dominates} over $\Cl{11}$ ($T_{1/2}=\SI{1223}{s}$). Since $\Ox{15}$ has
the longest positron range of the relevant emitters, the positron-range floor is
a genuine contributor (captured automatically by the PROD$\to$ANH displacement of
Sec.~\ref{sec:production}). In-beam and offline modalities are out of scope.

\paragraph{Simplification for a phantom.}
The paper's washout terms (fast/medium/slow tissue components) are \textbf{dropped}:
an inert PMMA or water phantom has no biologic clearance, so only physical decay
applies -- physically exact for our target.

\begin{table}[h]
\centering
\caption{Beam time structures (Parodi et al., their Table~3). The tabulated
$t_{\mathrm{del}}$ and $t_{\mathrm{meas}}$ are the paper's offline values; for the
in-room scope we use a short delay $t_{\mathrm{del}}\gtrsim\SI{120}{s}$, keeping
the accelerator $t_s,t_p,t_{\mathrm{irr}}$ columns.}
\label{tab:time}
\begin{tabular}{lcccccc}
\toprule
Accelerator & $t_s$ (s) & $t_p$ (s) & $t_{\mathrm{irr}}$ (s) & $t_{\mathrm{end}}$ (s) & $t_{\mathrm{del}}$ (s) & $t_{\mathrm{meas}}$ (s)\\
\midrule
Cyclotron  & 60  & 0   & 60  & 40 & 300 & 1800 \\
Synch.\ GSI & 1   & 4   & 600 & 40 & 300 & 1800 \\
Synch.\ HIT & 1.5 & 1.5 & 180 & 60 & 300 & 1800 \\
\bottomrule
\end{tabular}
\end{table}

\paragraph{The A$\to$B handoff.}
This is where stage~A meets stage~B, and it is the only place the bookkeeping
touches the simulation. For the chosen scenario, Eq.~\eqref{eq:off} gives $N_j$ for
each species; we then \textbf{draw $N_j$ ANH events from species $j$'s list} in the
PROD/ANH file (with replacement -- the list encodes the \emph{shape}
$P_j(\mathbf{r})$ at high statistics). Because $N_j$ is sampled per isotope, the
measured $\Ox{15}/\Cl{11}$ mixture comes out correct automatically, without any
reweighting. Drawing $N_j$ as a Poisson variate across $M$ realisations supplies
the noise for the $\sigma(\text{range})$ metric of Sec.~\ref{sec:compare}. The
result is the \textbf{sampled annihilation events} fed to stage~B; from each ANH
point, two \SI{511}{keV} photons are emitted back-to-back with acollinearity
($\sim\!0.5^\circ$ FWHM).

\section{Geant4 stage B: detector --- from annihilation events to the coincidence list}
\label{sec:detector}

This stage runs \textbf{once per candidate PET}. It reads the \emph{same} sampled
annihilation events and produces a \textbf{coincidence list} -- the deliverable of
the simulation and the sole input to reconstruction (Sec.~\ref{sec:recon}):
\begin{enumerate}[nosep]
\item \textbf{Detector simulation (Geant4):} the two \SI{511}{keV} photons
propagate through the phantom (attenuation, Compton scatter) and into the
detector; model intercrystal scatter, depth-of-interaction, energy resolution
$\sigma_E$ and time resolution $\sigma_t$.
\item \textbf{Coincidence sorting:} energy window, coincidence time window;
classify true/scatter/random. In the in-room scenario the in-beam
prompt-$\gamma$/neutron background is absent; the relevant backgrounds are random
and patient-scattered coincidences, plus -- for a LYSO candidate -- the intrinsic
${}^{176}\mathrm{Lu}$ singles floor, injected as a volumetric source in the
crystals (Sec.~\ref{sec:compare}).
\item \textbf{Coincidence list (the output).} Each accepted coincidence is written
with everything a list-mode reconstructor needs:
\begin{center}
the two \textbf{DOI-resolved 3-D hit positions} (or crystal IDs + a known geometry),
the two \textbf{energies}, and the two \textbf{timestamps} (for TOF).
\end{center}
Truth flags (true/scatter/random) may be appended for our own analysis; the
reconstructor does not require them.
\end{enumerate}

\section{Reconstruction (independent stage)}
\label{sec:recon}

Reconstruction is deliberately \textbf{decoupled} from the simulation: it consumes
only the coincidence list, so the toolkit (a list-mode MLEM/OSEM of our own, or an
external package such as CASToR or STIR) can be chosen or replaced later without
touching stages~A or~B. For the present plan it suffices to fix the requirements:
\begin{itemize}[nosep]
\item \textbf{List-mode} (not sinogram-binned): the measurement is photon-starved
and time-resolved (per-frame isotope separation), and the events carry per-event
DOI.
\item \textbf{DOI in the system model:} oblique lines of response are exactly where
it prevents parallax blur, so the reconstructor must use the 3-D hit positions, not
a fixed crystal-face radius.
\item \textbf{TOF-aware} when the timing axis is being compared.
\end{itemize}
The output is the 3-D activity image and, in particular, the 1-D depth profile
along the beam axis. Reconstruction must always be performed (not merely a
projection-space comparison), since the resolution, scatter-rejection and parallax
differences between detectors only fully manifest in the image.

\section{Comparing the detectors}
\label{sec:compare}

\paragraph{Same source, different responses.}
Every candidate PET is fed the identical annihilation events. Any difference in
the reconstructed images therefore comes from the detector (geometry, sensitivity,
$\sigma_E$, $\sigma_t$, DOI), not from the source. This is what makes the Geant4
cross-section weakness irrelevant to the \emph{ranking}: a systematic error in
$P_j(\mathbf{r})$ is common to all detectors and cancels. The file can later be
regenerated from a better production map without re-running any detector
simulation.

\paragraph{Figure of merit: range precision.}
The clinically meaningful quantity is the position of the activity distal
fall-off, which tracks the proton range (it sits a few mm \emph{proximal} to the
Bragg peak, because production needs above-threshold protons). For each detector:
\begin{enumerate}[nosep]
\item fix the count budget from Eq.~\eqref{eq:off} at $t_{\mathrm{del}}\approx\SI{120}{s}$;
for a single \SI{1}{Gy} field this is the photon-starved regime of $10^{7}$--$10^{8}$
decays, i.e.\ only $\sim\!10^{5}$--$10^{6}$ collected coincidences after attenuation
and efficiency;
\item generate $M$ independent Poisson realisations at that budget;
\item reconstruct each, fit the distal fall-off, and report the spread
$\sigma(\text{range})$ in mm.
\end{enumerate}
The headline plot is $\sigma(\text{range})$ \textbf{vs.\ counts} (equivalently vs.\
dose, or vs.\ delay $t_{\mathrm{del}}$, which sets how much $\Ox{15}$ survives) per
detector. Secondary metrics: recovered axial/lateral resolution, scatter fraction
retained inside the energy window, and parallax blur at large axial angle.

\paragraph{What to keep fixed vs.\ vary.}
Fixed across detectors: the annihilation-event source, the phantom, the count
budget, and the reconstruction algorithm family. Varied: only the detector. The
discriminators that matter for in-room range verification of this photon-starved,
scatter-dominated, $\Ox{15}$-rich signal -- and which the matrix is built to
exercise -- are:
\begin{itemize}[nosep]
\item \textbf{Solid-angle sensitivity (AFOV).} The leading requirement. Sweep the
axial length $L$; the geometric acceptance for a centred source scales roughly as
$L/\sqrt{L^2+R^2}$, the dominant lever on how many of the scarce pairs are
collected.
\item \textbf{Energy resolution $\sigma_E$.} Sets how tightly the \SI{511}{keV}
window can be drawn, hence how much patient scatter is rejected on the
\emph{oblique} coincidences a long AFOV adds. Compare e.g.\ $\sim\!6\%$ (cryogenic
CsI) vs.\ $\sim\!10\%$ (LYSO) vs.\ poor (BGO).
\item \textbf{Depth-of-interaction.} Controls parallax on those same oblique LORs.
Compare DOI-capable monolithic crystals against fixed-depth pixels.
\item \textbf{Intrinsic radioactivity.} Inject the ${}^{176}\mathrm{Lu}$ singles
floor for the LYSO arm (none for CsI/BGO); at sub-MBq signal it raises randoms and
competes with the signal, so it is part of the comparison, not a detail.
\item \textbf{Timing / TOF.} Treated as a \emph{low-value} axis here: include a
TOF-on variant to demonstrate diminishing returns for a one-dimensional
distal-edge task, rather than as a discriminator the candidate relies on.
\end{itemize}
\textbf{Candidate set:} cryogenic monolithic CsI (the CRYSP baseline) against
conventional LYSO and BGO scanners, plus the CRYSP variants -- room-temperature
CsI(Tl), cryogenic BGO, and the BGO-core/CsI-wing hybrid -- each represented by its
$(\text{geometry},\,\sigma_E,\,\text{DOI resolution},\,\text{intrinsic background})$.
Since all see the identical source, the resulting $\sigma(\text{range})$ vs.\
counts curves isolate exactly these detector properties.

\section{Scope and honest caveats}
\label{sec:scope}

\begin{itemize}[nosep]
\item \textbf{Homogeneous phantom} omits tissue heterogeneity, hence no
$\Ox{15}/\Cl{11}$ spatial split and no washout (deliberate; this is a detector
study). A two-region phantom (e.g.\ a bone-like insert) can be added later to
probe density-interface artifacts.
\item \textbf{Geant4 cross sections} bias absolute counts and the production
shape; the relative comparison is robust, the absolute range precision is
indicative.
\item \textbf{In-room scope.} Only the in-room modality is treated (full-ring
scanner, $t_{\mathrm{del}}\gtrsim\SI{120}{s}$); in-beam and offline are out of
scope. At this short delay the source is $\Ox{15}$-dominated.
\item \textbf{Positron range} ($\Ox{15}$, $\Cl{10}$) sets a physical resolution
floor; a high-performance PET may saturate against it -- a result to report, not
to hide. It is captured by the PROD$\to$ANH displacement.
\item \textbf{Discarded prompt $\gamma$ / secondary neutrons} (Sec.~\ref{sec:production})
are small omissions, acceptable for a first cut and listed here for honesty.
\end{itemize}

\end{document}
